\subsection{Overview}

Because the location of buried tile networks is frequently undocumented, poorly documented, or out of date, a substantial body of work has established that non-invasive remote sensing and near-surface geophysical methods can detect and map agricultural subsurface drainage. The literature spans two complementary families of technique: (i) aerial and unmanned aerial vehicle (UAV) imagery in the visible (VIS/RGB), multispectral (MS), near-infrared (NIR) and thermal infrared (TIR) bands, which senses drainage indirectly through soil-surface and vegetation contrasts, and (ii) near-surface geophysical methods---ground-penetrating radar (GPR), electromagnetic induction (EMI), electrical resistivity tomography (ERT) and time-domain reflectometry (TDR)---which sense the pipes or the moisture/resistivity anomalies around them more directly. Systematic reviews spanning dozens of studies conclude that no single method or spectral band is universally superior, and that detection success depends strongly on soil moisture, soil composition, drain depth and diameter, and the timing of surveys relative to rainfall, irrigation and field operations \parencite{Carlsen2024,Wienken2024}. The strongest consensus is that combined strategies---particularly UAV RGB/TIR imagery together with GPR---are the most promising, and that timing and antecedent moisture are decisive controls on detectability \parencite{Wienken2024,Carlsen2024,Koganti2021}. Complementing detection, a further line of evidence documents GIS/GPS and digital-elevation-model (DEM) based tools for planning, designing and managing drainage systems once their layout is known \parencite{Shedekar2018,Falasy2024,Falasy2025,Falasy2025b}. This section synthesises what these lines of evidence have established.

\subsection{UAV and aerial imagery: spectral bands and detection rates}

Multiple field campaigns have quantified how well VIS/RGB, MS and TIR imagery reveal drain lines through surface soil-moisture and vegetation patterns. In the largest multi-site assessment, spanning 29 sites across Indiana, Iowa, Michigan and Ohio, UAV visible-color (VIS-C) imagery detected at least some drain lines at 48\% of sites (14 of 29), multispectral imagery at 59\% (17 of 29), and thermal infrared imagery at 69\% (20 of 29); TIR generally performed best overall, although at some sites VIS-C or MS provided better insight into drain-line patterns than TIR \parencite{Allred2020}. The authors conclude that all three imagery types have significant potential for mapping drainage pipe systems, that the best practice at any site is to acquire all three, and that even where drainage detection fails the imagery remains useful for the topographic information it provides \parencite{Allred2020}.

\subsubsection{Thermal infrared imagery}

TIR has emerged as the most consistently effective single band under many conditions. In Ohio, UAS TIR imagery detected roughly 60\% of the known subsurface drainage infrastructure, whereas under very dry conditions drainage pipes were not detected with VIS and NIR imagery; the authors conclude that TIR has substantial promise and may be better suited than VIS or NIR under dry surface conditions \parencite{Allred2018}. A downscaled Illinois experiment established the feasibility of using thermal rather than optical imagery, showing that raw optical images reveal no detectable tile-related features and that optical-based processing fails, while thermal images provide higher contrast for distinguishing regions near and away from tiles; the proposed thermal processing approach successfully detected and localised pipes even in the presence of tillage-induced surface line features \parencite{Woo2019}. That work also found no strong relationship between soil and atmospheric temperature, supporting the hypothesis that surface soil-temperature differences are driven by tile-related moisture patterns, and identified the optimum thermal acquisition window as within 12~h after 2.54~mm of precipitation \parencite{Woo2019}. The method was presented as an effective, promising tool for precisely locating tile systems over large areas \parencite{Woo2019}.

The magnitude of the exploited thermal signal is small: UAV TIR surveys in Ohio measured apparent radiant temperature differences of less than 1.0\,\textdegree C between soil over drain lines and soil between them, yet this was sufficient to map subsurface drainage patterns \parencite{Allred2021}. These surveys clearly showed substantial promise for mapping drainage in farm fields \parencite{Allred2021}. Timing within the day mattered: late morning through late afternoon was identified as the best time based strictly on consistency of success, whereas early morning and early evening were transition periods when no TIR drain-line response may be observed and should be avoided; near sunrise or sunset, drain lines can appear as darker shaded linear features, often easier to distinguish from lighter farm-operation features \parencite{Allred2021}. Meteorological conditions also constrained data quality: high relative humidity (>60\%) can prevent production of a complete orthomosaic by degrading TIR image quality, the relationship between percent calibrated images and relative humidity was negative, air temperature and dew point showed very low correlations with orthomosaic production, and the positive wind-speed correlation was described as misleading \parencite{Allred2021}.

TIR has also been applied to locate drainage outlets in surface water. In a preliminary feasibility assessment across five sites, UAV TIR detected outlet-related thermal responses at a small stream (Site~1) and a drainage ditch (Site~2), with about a 1\,\textdegree C change in water apparent radiant temperature from upgradient to downgradient of the outlet in each case, but did not detect responses at the three moderate-to-large streams (Sites~3--5) \parencite{Allred2022}. At Site~1 the direction of the thermal response reversed between cold winter conditions and warmer periods of the year \parencite{Allred2022}. The authors caution that other factors in ditches and streams may produce water temperature variations of similar magnitude to a pipe-outlet response, limiting practical outlet mapping in some settings \parencite{Allred2022}.

Thermal detection is not universally successful. At Illinois sites, thermal imagery showed only limited evidence of thermal contrast related to the drainage pipe and did not provide sufficient contrast for detection, with the thermal signature falling below the sensor's thermal resolution; the low contrast was attributed to a modest preceding rain event, the long interval between rain and acquisition, or persistent hot wind during the survey \parencite{Kratt2020}. In that study RGB imagery clearly revealed the pipe network at one site and provided the best guide to where pipes might be detected in the thermal data, and the timing of survey collection relative to the precipitation event was emphasised as important \parencite{Kratt2020}. Both thermal and RGB imagery nonetheless demonstrated potential to elucidate drainage pipes, and UAS thermal and RGB surveys were characterised as a transformational method for agricultural water management \parencite{Kratt2020}.

\subsubsection{Visible, multispectral and photographic imagery}

Beyond TIR, visible and multispectral imagery contribute independent information, and combining bands improves completeness. Some combination of VIS-C, MS and TIR can provide more complete drainage-pattern information than any single type alone \parencite{Allred2020}. Longer-standing aerial photographic work in the Czech Republic showed that functional drainage lines could be recorded both during and outside the vegetated period owing to visual differences in soil surface or vegetation, and that the primary drainage-feature type captured in panchromatic natural-color images correlated with that identified in infrared photos, with no differences in basic drainage type or presence/absence between synchronously obtained images \parencite{Tlapáková2017}. Importantly, all test localities revealed drainage systems not registered in documentation, older systems were often missing from the registry, and differences between actual and plotted lines confirmed the need to revise the national amelioration database \parencite{Tlapáková2017}. Field verification showed that drain condition and functionality, drainage-line characteristics, building technology and current soil condition controlled surface manifestations such as local waterlogging, and success varied across sites with meteorological and spatial conditions during survey flights, with only a small fraction of identified systems falling into the "complete and well visible" category \parencite{Tlapáková2017}. The method was judged an effective means of acquiring up-to-date, precise digital documentation of drainage position and functionality; the authors recommended combining standard large-scale aerial imaging with complementary UAS recording of smaller parcels, combining imagery with information on actual field conditions, and noted that detailed digital surface models and more frequent imaging may further improve mapping \parencite{Tlapáková2017}.

\subsection{Ground-penetrating radar and antenna configuration}

GPR directly senses buried pipe and constructed-layer reflections and has been demonstrated across a range of settings, with performance strongly modulated by soil conditions, antenna orientation and survey geometry. Early work on golf-course tees and greens found that antenna center frequencies of 250--1000~MHz worked reasonably well for mapping buried drainage pipe networks and delineating the areal extent of constructed soil layers, while higher frequencies (especially 900--1000~MHz) were best for determining depths and thicknesses of constructed layers; computer-modelled synthetic GPR profiles of different design scenarios guided equipment selection and processing and aided interpretation of field profiles---for example in picking the bottom of constructed sand and gravel layers and identifying reflection hyperbolas corresponding to buried pipe \parencite{Allred2005}. EMI, using a ground conductivity meter with relatively short intercoil spacing at the surface, was responsive to constructed sand/gravel layers, and the study concluded that both EMI and GPR provide valuable information for tee and green maintenance and remodelling \parencite{Allred2005}.

Antenna orientation relative to drain lines interacts strongly with soil moisture. In an Ohio case study, under moderately dry conditions (August~12) antennas perpendicular to the drain lines produced the strongest responses, with scaled, signal-enhanced maximum pipe-response amplitudes averaging 50\% greater than the parallel orientation across Pipes~1--4, A and B (and 120\% greater for the Pipe~2 banded response); under saturated wet conditions (August~13) the parallel orientation was strongest, with overall amplitudes about 25\% greater than perpendicular (15\% greater for Pipe~2); and under drained wet conditions (August~14) parallel orientation was marginally stronger (about 10\% greater for Pipe~2) \parencite{Allred2013}. These results provide practical guidance on GPR antenna set-up for locating and evaluating field drainage pipes \parencite{Allred2013}.

GPR integrated with RTK/GNSS positioning has been shown to be a valuable, efficient tool for determining drainage-pipe patterns. Across three agricultural sites in Maryland and Ohio it successfully delineated pipe patterns: a complex rectangular system at Site~1 (southwest--northeast lines on the west, southeast--northwest on the east), enough of a herringbone system at Site~2 to unambiguously identify its type despite incomplete detection, and random, likely very old lines at Site~3 interpreted as draining wet areas of a hay field \parencite{Allred2018b}. Soil homogeneity aided detection---responses were easier to isolate in the more homogeneous soils at Site~1 than in more heterogeneous soils---and a novel use of spiral or serpentine transect segments established drain-line presence and directional trends at Sites~2 and 3 \parencite{Allred2018b}. Data collection was efficient, requiring only four to six hours per site without a dense transect grid, and the approach was presented as a substantial improvement over tile-probe and trench-excavation methods because it was faster, better able to resolve drain-line trends, and caused no pipe damage \parencite{Allred2018b}.

Three-dimensional GPR extends these capabilities but remains soil-limited. Across twelve Danish sites, a high success rate in mapping drain pipes with 3D-GPR was achieved at five sites while the other seven were less successful; detection was generally better when the survey direction was perpendicular or at an angle to the expected drain orientation, and surveys parallel to the drain lines reduced detectability \parencite{Koganti2020}. The inverse relationship between 3D-GPR penetration depth and mean electrical conductivity (0--1.5~m) was generally confirmed, supporting the idea that higher soil EC increases signal attenuation, although localised penetration depth and EC showed site-specific behaviour with strong-to-moderate negative correlations at some sites and weak or no correlation at others \parencite{Koganti2020}.

\subsection{Geoelectrical and injection-based methods}

Time-lapse ERT, often coupled with an injected contrast agent, has been shown to locate drainage networks by exploiting induced resistivity changes. On a small French agricultural plot, field time-lapse ERT coupled with water injection located the drainage network, with drains clearly defined after sufficient injected water; ERT and TDR measurements showed similar temporal trends, with deeper horizons affected earlier than surface layers, and their good correlation was said to confirm and validate ERT for this application \parencite{Jouen2018}. Synthetic 3D modelling delimited the method's sensitivity: infiltrations 0.3~m thick were not detected regardless of width (0.4--1.0~m) or bulk resistivity ratio (0.89--0.69), whereas detection was best for thicknesses greater than 0.55~m and widths greater than 0.7~m \parencite{Jouen2018}. The authors report this as the first experimental demonstration of preferential flows beneath a drainage system, conclude that water injection combined with ERT is a good solution for locating drainage systems, but note that 3D time-lapse ERT reaches its limits on larger agricultural plots \parencite{Jouen2018}.

An alternative contrast mechanism uses pressurised air rather than water. Laboratory and field experiments demonstrated that coupling ERT monitoring with high-pressure air injection from the drainage outlet detects drain position in both space and depth, with high-pressure injection causing an instantaneous increase in electrical resistivity around the drains \parencite{Henine2021}. In the laboratory, 10~mbar produced only small resistivity changes, whereas 30, 39 and 55~mbar produced significant increases; in the field, low-pressure injection at 50~mbar with warm air (25--35\,\textdegree C) did not increase resistivity, while high-pressure injection at 100 and 150~mbar produced a significant resistivity contrast around the drain after about 10~min \parencite{Henine2021}. High-temperature air injection did not raise soil temperature around the drains enough to affect the resistivity signal, indicating the response was driven by displacement rather than heating \parencite{Henine2021}. The previously known drain positions were successfully located, and the authors present the method as offering new perspectives for detecting buried pipes \parencite{Henine2021}.

\subsection{Complementary and multi-method strategies}

A recurring, well-supported finding is that pairing methods with different sensitivities yields more complete and better-validated maps than any single technique. Across four Midwest USA sites (Ohio and Michigan), UAV imagery and GPR were complementary for mapping subsurface drainage \parencite{Koganti2021}. At Site~1 both detected drainage patterns and GPR additionally confirmed pipe locations and provided depth where UAV imagery faded; at Site~2 recent UAV imagery under cover crop revealed drainage locations while GPR was useful where UAV imagery showed no pipe signature; at Site~3 UAV imagery on bare ground revealed no drainage signatures but GPR confirmed the field was drained and determined pipe orientation, with later UAV imagery under crop cover showing only subtle differences; and at Site~4 UAV imagery was mostly successful after rainfall with better timing, whereas GPR had limited success owing to shallow penetration depth and noisy data at some locations \parencite{Koganti2021}. Overall, UAV imagery was described as feasible and cost-effective while GPR served as a complementary mapping and validation technique that also provided depth and orientation information, and the combined UAV--GPR approach was presented as more optimal, with guidance on optimal survey timing and a recommendation to use both technologies when available \parencite{Koganti2021}.

Systematic reviews reinforce and generalise these conclusions. A synthesis of 19 peer-reviewed remote sensing studies (15 focused on individual tiles, 4 on drained areas/fields) found visual inspection to be the most common approach, with geoprocessing and deep learning used for individual-tile detection only in the oldest study and in four of the most recent; imagery acquired on bare soil and after rainfall was generally most favourable, TIR often performed well, no single band was universally superior, and reported accuracies varied widely such that differing evaluation procedures make robust comparison difficult \parencite{Carlsen2024}. That review concludes that deep learning shows the best performance for large-scale delineation, that multitemporal acquisition may help reduce confusion with linear farm features, and it identifies the need for standardised evaluation frameworks and publicly available drainage-map datasets \parencite{Carlsen2024}. A second review screening 69 relevant publications found that no single geophysical or remote sensing method consistently demonstrated high success, and concluded that combining GPR and UAV technologies with RGB and TIR imagery is the most promising approach \parencite{Wienken2024}. It documented method-specific favourable conditions: for GPR, moderate soil moisture and lower-EC soils; for EMI, low soil moisture and higher-EC soils; for ERT, high soil moisture and electrically resistive soils; and for remote sensing, RGB best on bare soil or during vegetation, MS during active crop growth, and TIR in late morning to early afternoon under recent irrigation or rainfall \parencite{Wienken2024}. Detection success was emphasised to depend strongly on soil moisture, soil composition, drainage depth and diameter, and timing relative to agricultural and weather conditions \parencite{Wienken2024}.

A cross-cutting practical challenge is that linear features caused by drain lines can be confused with linear features produced by field operations, but knowledge of drainage-system installation and field operations helps distinguish them \parencite{Allred2020}. Timing relative to rainfall is repeatedly decisive: surveying too soon after rainfall sometimes prevented detection, while in other cases drain lines became clear only hours or a day later as soils drained and dried \parencite{Allred2020}.

\subsection{Large-scale and watershed delineation}

Beyond field-scale detection, machine learning and hydrologic filtering extend mapping to regional and watershed scales. High-resolution subsurface drainage maps have been developed using satellite big data and random forest machine learning via Google Earth Engine, motivated by the very limited availability of reliable drainage records in the United States despite their importance for sustainable water-resource management; in this framework soil variables were important for identifying potential drainage areas while land surface temperature distinguished where drainage had actually occurred \parencite{Cho2019}. At the watershed scale in Canada and Belgium, a remote sensing approach achieved 87.8\% accuracy in the first study area but only 17.3\% in the second, whereas a decision-tree classification identified potential drainage units better, capturing nearly 96.8\% of the ditch network in the first area and about 58\% in the second \parencite{Yimer2023}. Adding groundwater-level-based filtering improved identification by contributing 19.4~km\textsuperscript{2} of drained area not detected by remote sensing, and shallow groundwater-table maps closely followed the ditch network in the second area, supporting the approach \parencite{Yimer2023}. The study concludes that combining remote sensing, decision-tree classification and groundwater-table information---a novel complementary filtering step---improves and accurately identifies drainage units \parencite{Yimer2023}.

\subsection{GIS-based design, mapping and management tools}

Once layouts are detected or planned, GIS/GPS and DEM-based tools support design and management. GPS, remote sensing, DEMs and related software can integrate spatial and temporal information on soil type, hydrology and crop management at a site, enabling precise conservation practices and improved drainage management; overlaying georeferenced drainage layouts on DEMs or contour maps can delineate zones of water-table elevation for controlled drainage and subirrigation, and improved drainage management could raise yields and enhance soil and water conservation while helping adapt to a changing climate \parencite{Shedekar2018}. Recent advances allow frequent use of precision information for more effective, lower-impact drainage management \parencite{Shedekar2018}.

A suite of publicly accessible QGIS plugins has been developed and tested in Illinois to streamline drainage design. The tools ease the laborious process of identifying high points, plateaus, draws and depressions from contour lines and integrating them into a layout: a LiDARThAn tool reduces LiDAR elevation datasets to a specified pixel size and creates a dot at each pixel centre, while a Plot Field Laylines tool detects depressions by contrasting filled and unfilled DEMs to guide placement of mains and outlet points \parencite{Falasy2025}. Run consecutively, the four QGIS-based tools produced topologically sound, unidirectional flow tile networks and were judged efficient---executing objectives in simple "one clicks" and providing layout-design guidelines---with user testing finding them simple, user-friendly and much faster than the independent-task or direct methods \parencite{Falasy2025}. A companion toolset addresses topology and flow ordering: it identified and fixed positional digitisation errors (fractured lines, overshoots and dangles), a Network Flow-Path Generator produced a layout with zero unconnected segments, and QGIS Topology Checker found 71 errors in a layout created with commercial drainage software versus 34 in the layout from the new tools \parencite{Falasy2024}. Additional tools produced flow-ordering outputs and added elevation, length, slope and cumulative flow-length fields; the plugins were easy to use, reduced the expertise required, and their complete combination of techniques for identifying and fixing fractured line networks was highlighted as a major novelty \parencite{Falasy2024}.

Extending to depth placement and sizing, interactive QGIS tools automate and optimise drainpipe depth placement and estimate pipe sizes for each segment, sequencing both downstream-to-upstream and upstream-to-downstream \parencite{Falasy2025b}. Tested on a tile network, they produced exportable outputs for buried elevation depths, nominal pipe sizes and price estimations for the entire distribution; validation showed burying slopes consistent with existing software, although the new tools sometimes generated smaller sizes---for the outlet segment they gave an 11.15\,\textquotedbl\ actual and 12\,\textquotedbl\ nominal size versus 12.88\,\textquotedbl\ actual and 15\,\textquotedbl\ nominal from the existing software \parencite{Falasy2025b}. Critically, the new tools completed burying and sizing in under 10 minutes whereas the proprietary software required over two hours, and the authors state that the tools improve workflow efficiency, transparency and accessibility of subsurface drainage design within a GIS environment, with exportable outputs enabling cost analysis and integration with drainage-installation machines \parencite{Falasy2025b}.

\subsection{Drain envelopes and near-drain flow physics}

Finally, interpretation of geophysical and hydraulic signals near drains rests on an established understanding of drainpipe envelopes and near-drain soil processes. Thin envelope materials substantially reduce entrance resistance, with the reduction significant up to about 5~mm thickness and up to a ratio of envelope-to-soil hydraulic conductivity of 20; for less permeable surrounds, entrance resistance increases substantially up to about 10~mm thickness and for conductivity ratios below 0.2, while at conductivity ratios of 10 and above envelopes largely eliminate differences in entrance resistance between drain pipes \parencite{Stuyt2006}. Approach-flow resistance is lowest for fully radial approach flow or conditions engaging the entire drain circumference, and it decreases with increasing discharge toward a constant minimum \parencite{Stuyt2006}. Field observations and CAT-scan studies showed that flow resistance near a drain is largely associated with soil macroporosity and the geometric arrangement of the near-drain macropore network, and that good installation practice favours the service life of wrapped drains whereas well-designed envelopes cannot compensate for unfavourable soil properties or poor installation \parencite{Stuyt2006}. These features of drains, envelope materials and near-drain physical processes underpin the moisture and resistivity contrasts that remote sensing and geoelectrical methods ultimately exploit \parencite{Stuyt2006}.